% =====================================================================
%  R. Nielsen, FILOSOFI OG SUND MENNESKEFORSTAND
%  In: For Ide og Virkelighed. Et Tidsskrift, udgivet af R. Nielsen, B. Bjørnson og Rud. Schmidt.  Aargang 1870. Første Bind (Bind 2).
%  Kjøbenhavn: Chr. Steen & Søns Forlag; Louis Kleins Bogtrykkeri.  Bind 2, pp. 385-404.  Antiqua.
%
%  TRANSCRIPTION -- GENERATED; DO NOT EDIT BY HAND.  Rasmus Nielsen's eight essays in the volumes 1870-1873 are built by ONE
%  pipeline (ebook-method/for-idee-1870-73/) into a combined working file, which split7.py cuts into seven transcriptions
%  (Aandsvirkelighed and its sequel, Videre om Aandsvirkelighed, are one: Hans, 2026-10-13).
%  Rebuild from the repo root:   sh ebook-method/for-idee-1870-73/rebuild.sh
%
%  STATUS: COMPLETE (2026-10-13): the volume's title page and the whole essay, Bind 2, pp. 385-404, with its notes.
%
%  EDITION.  The only printing.  Text of record: Nasjonalbiblioteket's scan (nb.no) of the volume, ~/bibliotek/Nielsen, Rasmus/1870-1-for-idee-og-virkelighed.pdf (NB item 440261bfdbcab019d4f313aa5894cf6e) (SCANS.tsv),
%  page images fetched by nbfetch.py, 300 dpi colour, NO text layer; printed page N = PDF page N + 4.  No e-book, no
%  Wikisource text, no reprint, no earlier transcription found (checked 2026-10-13: the repo, the web, the Internet Archive,
%  da.wikisource; HathiTrust and the Google Books API could not be queried that day; Google Books holds at least the 1873
%  volume, nielsen/darwinismen's source -- a second copy, not used).
%
%  ACCURACY (2026-10-13; one collation for all eight texts, 231 pp., ~58,000 words).  The DRAFT is NB's own OCR (ALTO, ABBYY
%  FineReader 12), laid over the page images (combine.py) and turned into LaTeX by draft.py; it was collated word by word
%  with an independent tesseract reading of the same images (collate/check.py, witness 'ocr'): 1556 disagreements.
%  check.py's own rules settled 78.  Candidate rules were first scored BY SCRIPT on a text already read at the image
%  (calibrate.py: Et Synspunkt for Darwinismen, same volume as Verdensanskuelse), then measured on a BLIND PILOT of 152 crops
%  (reader A, 52/52 planted controls right); 17 rules passed (NB's glyph faults: o for ø, D/O/U for Ø, i for l, li for h,
%  l for I, opening » read «, words run together or broken apart; tesseract's: aa read an/ag/å, added accents; the
%  printer's space before a mark, house rule) and settled 772; the rest were read at the image BLIND (562 crops on 94
%  sheets, four readers, 188/188 planted controls right), ~160 decided in the main session from the readers' lines and ~40
%  at 300-dpi zooms.  STRUCTURE read at the image (reader K): every footnote whole (30; collated where they stand, then moved
%  to their marks, then compared word by word with the reading: notecmp.py, 9 differences settled at the image), the
%  headings, the section numerals, the essays' last lines and signatures, the three Greek words, the volume title pages.
%  LETTERSPACING (the print's emphasis, \emph): candidate lines by letter gaps and box width, calibrated on the words the
%  word readers had marked spaced (spaced.py): 361 lines read blind (61/62 roman controls), 25 doubtful extents read at a
%  zoom.  SWEEPS after OPEN 0: words run together that both OCRs share (lower-case + capital), the unknown-word sweep (678
%  words, ~20 suspects, 14 read at the image: 4 OCR faults, 7 misprints kept), --moved, the dropped-word list, the
%  W-shortens list of the blind decisions.
%  SAMPLED PAGE CHECK (pagecheck.py, seed 202610137): one page of each essay or part, pp. 2289, 2400, 3305, 4308, 5033,
%  5208, 6514, 8011 of the working numbering (Bind x 1000 + page), read whole and blind, 1965 words: 0 word errors,
%  1 punctuation (a compound's hyphen lost at a page turn: „Ny-|Holland“) and 1 letterspacing („to be or not to be“)
%  error, both mended.  95 % upper bound for word errors: 1.5 per 1000 words.
%  What this cannot see: a misprint both OCRs normalise into a real word.
%  THE TRANSLATION AS A SECOND READING (2026-10-14): the 15 translators listed 74 doubts; ~50 were read at 300-dpi
%  zooms: 13 faults mended (4 compound hyphens set with a space, 4 specks taken for marks, a lost „I“, a colon read as
%  a semicolon, a comma read as a stop, „à“ read „å“, a space before a mark) and 9 false paragraph breaks at page
%  turns removed (left by the notes moved away; ALTO shows the next page flush); 3 more misprints kept (PRINTED AS IS).
%  The ø rule (Hans 2026-10-08): o-forms that are no word are set with ø; „stode“ (= stood) is a word and stays.
%  PRINTED AS IS Bind 2, p. 392: „Naturhistorie o, s. v.“ comma for the stop of o. (b1-19-04)
%  PRINTED AS IS Bind 2, p. 389: „Livggaaderne“ for Livsgaaderne (zoom z4)
%  Not transcribed: running heads, folios, sheet signatures („For Ide og Virkelighed. 1870. I. 20“), library stamps; the
%  large initial of each essay is set as a plain letter.  Quotation marks as printed (» « in the text, sometimes a closing
%  »; „ “ in some notes).  The notes are numbered ¹), ²) per page, as printed.  Greek as printed (textalpha).
%
%  TOOLS (ebook-method/for-idee-1870-73/): nbalto.py, combine.py, alto.py, ocr.py, draft.py, rules.py, calibrate.py,
%  sheets.py + reads.py + READER-BRIEF.md, spaced.py + spaced.tsv, structsheets.py, notecmp.py, unknown.py, zoom.py,
%  pagecheck.py + PAGECHECK-BRIEF.md, patches.py (PATCHES, PRENOTES, NOTES, POST, MISPRINTS), decisions.tsv, split7.py,
%  rebuild.sh.
% =====================================================================
\documentclass[12pt,a4paper,openany]{book}
\usepackage[utf8]{inputenc}
\usepackage[T1]{fontenc}
\usepackage{libertinus}
\usepackage{libertinust1math}
\usepackage{textalpha}
\usepackage[danish]{babel}
\usepackage[protrusion=true,expansion=true]{microtype}
\usepackage{setspace}
\usepackage[margin=1.3in]{geometry}
\usepackage{fancyhdr}
\usepackage{hyperref}
\hypersetup{pdftitle={Filosofi og sund Menneskeforstand}, pdfauthor={Rasmus Nielsen}, hidelinks}
\pagestyle{fancy}
\fancyhf{}
\fancyhead[LE,RO]{\thepage}
\fancyhead[C]{\textit{R. Nielsen: Filosofi og sund Menneskeforstand (1870)}}
\setstretch{1.3}
\usepackage{marginnote}
\renewcommand{\thefootnote}{\arabic{footnote})}   % the print numbers each page's notes ¹), ²)
\newcommand{\opage}[1]{\marginnote{\footnotesize\textit{[#1]}}}
% the essay's heading: title (and subtitle), centred, a short rule; the print's large initial is not reproduced
\newcommand{\essayhead}[2]{\begin{center}{\large #1}\ifx&#2&\else\\[0.3em]{\normalsize #2}\fi\\[0.5em]\rule{3em}{0.4pt}\end{center}\medskip\noindent}
\newcommand{\signature}[1]{\par\smallskip\hfill #1\hspace*{2em}\par}

\begin{document}

% --- front matter: the title page of Bind 2 (PDF 1 of 1870-1-for-idee-og-virkelighed.pdf), read at the image (reader K, m1);
% short wavy rules are set as plain rules; the library stamps are not transcribed.
\begin{titlepage}
\begin{center}
\vspace*{3em}
{\Huge For Ide og Virkelighed.}\\[1em]
\rule{8em}{0.4pt}\\[1em]
{\large Et Tidsskrift,}\\[0.8em]
{\small udgivet}\\[0.4em]
{\footnotesize af}\\[0.6em]
{\large R. Nielsen,}\\[0.4em]
{\large B. Bjørnson og Rud. Schmidt.}\\[1em]
\rule{8em}{0.4pt}\\[0.6em]
{\large Aargang 1870. Første Bind.}\\[0.6em]
\rule{8em}{0.4pt}\\[4em]
Kjøbenhavn.\\[0.4em]
Chr. Steen \& Søns Forlag.\\[1.5em]
{\small Louis Kleins Bogtrykkeri.}
\end{center}
\end{titlepage}

% ===== TEXT: pp. 385-404 =====
\essayhead{Filosofi og sund Menneskeforstand.}{}
% --- p. 385 ---
\setcounter{footnote}{0}%
\opage{385}Det er endnu kun lidt over tredive Aar siden, at J.
L. Heiberg fandt sig i vor Literatur opfordret til at
tage Ordet for en dybere gaaende filosofisk Tænkning
og med Eftertryk minde om, hvor langt vi i denne
Henseende stode tilbage for andre civiliserede Nationer,
eftersom Filosofien endnu kun var en fremmed, neppe
engang velkommen Gjæst her i Landet. »Man har«,
siger han, »allerede ofte fremsat den Bemærkning, at
medens vi Danske tilegne os den almindelige europæiske
Kultur i dens mange forskjellige Retninger, er der dog
endnu een tilbage, ad hvilken vi hidtil ikke have ladet
os føre, nemlig den filosofiske. Det politiske, det videnskabelige,
det æsthetiske og det religiøse Liv bærer
hos os Nutidens almindelige Præg; fra alle hine Hjørner
blæser Kulturens velgjørende Vind paa vort lille Danmark,
som dog ikke blot ligefrem følger den herskende
Luftstrøm, men ogsaa mangen Gang bruger Seil og
Ror til at modificere sin Retning, til at benytte det
mægtige Element, uden at lade sig beherske af det.
Men idet vi villigen lade os gjennemblæse af den \emph{Aande},
der ifølge den nordiske Fabel kommer fra de fire Dværge
% --- p. 386 ---
\setcounter{footnote}{0}%
\opage{386}i Verdenshjørnerne af Ymers uhyre Hjerneskal, frygte
vi derimod for at forkjøles af den universelle \emph{Aand},
og krybe forsigtigen i vort Sneglehus, hvergang vi spore
dens Nærværelse« (Perseus. Nr. 1. 1837).

Meget Vand er siden den Tid løbet i Stranden.
Rigets Forfatning, Folkets Tanker, Literaturens Udseende,
det aandelige Livs Retning, Alt har i et Mellemrum af
faa Aar undergaaet en saadan Forandring, at det Gamle
og det Nye allerede staae som to Rethavende, der iagttage
hinanden mistænksomt, maale hinanden med Øinene,
og selv, hvor der paa begge Sider er god Villie, have
Vanskelighed nok ved at forstaae hinanden. Under de
stærke Bevægelser i det Udvortes er Blikket tillige blevet
vendt mod det Indvortes; en stigende Trang til aandelig
Selvbesindelse yttrer sig, om end paa forskjellig
Maade, i stedse videre Kredse. Saaledes er da Filosofien
hos os ikke længer en overflødig Luxus, en indført
Zirplante, der kun kan trives i et dertil opvarmet
Drivhus, men en uafviselig Fornødenhed, et Aandens
Krav, vi paa vore Vilkaar maae see at tilfredsstille,
saa godt vi kunne.

Thi hvad er Filosofi? En videnskabelig Bevidsthedslære,
en Videnskab, der vil opklare, hvad det er,
vi kalde vort Selv, og forhjælpe det tænkende Menneske
til at orientere sig i sin egen Bevidsthed. Alle
Slags Indbildninger og Sandsebedrag, Skuffelser og Selvbedrag
af enhver Art have, menneskelig talt, deres Rod
i Uvidenhed om det selvbevidste Livs Grundlove. At
ride vild i sin egen Skov er farligt, skjøndt i en vis
Forstand herskabeligt, forsaavidt det forudsætter, at
Skoven og dermed Eiendommen maa have en vid Ud%
% --- p. 387 ---
\setcounter{footnote}{0}%
\opage{387}strækning; men at løbe vild i sin egen Bevidsthed og
spørge som Bonden: »var det mig, der var her før?« er
jammerligt og det i samme Forhold som de individuelle
Naturgaver ere betydelige, og de aandelige Besiddelser,
Kundskaber, Erfaringer o. s. v., udstrakte. Der er da,
som man seer, Noget i Filosofien, der vedkommer Alle.
Ingen kan tænke over sit Liv uden at filosofere; Virkeligheden,
det daglige Liv med Omstændighederne og
de uforudseelige Tilfælde har Problemer nok: Jeppe i
Baronens Seng er jo nær ved at blive kvalt af Problemer.

Imod de nationale Bestræbelser i Nutiden har man
nu og da villet opstille den gamle Sætning, at Sandheden
er ens for Alle, at folkelige Forskjelligheder ikke
tør udøve nogen Indflydelse paa Videnskaberne, at navnlig
Filosofien, den almene Videnskab, maa være kosmopolitisk.
Forsaavidt her endnu tænkes paa Tænkningens
rene Love, har Sætningen ganske vist sin Gyldighed.
Men naar Spørgsmaalet er om den eiendommelige Behandling
af de videnskabelige Opgaver og om den ved
Folkesprogene betingede Fremstilling af Videnskabens
Indhold, gjør Nationaliteten sin Ret gjældende. Jo oprindeligere
de videnskabelige Frembringelser ere, og jo
mere Undersøgelserne nærme sig det inderste Aandelige,
desto stærkere vil Erkjendelseslæren og dermed den hele
Literatur være præget af Folke-Aanden. At Filosofien,
den kosmopolitiske Videnskab, ingenlunde kan være
tjent med at rykkes op af den folkelige Jordbund og
svæve frit i almene Tankers luftige Regioner, lærer dens
egen Historie. Grækernes Filosofi er en Afspeiling af
den græske Aand, og det er ingenlunde tilfældigt, at
% --- p. 388 ---
\setcounter{footnote}{0}%
\opage{388}Filosofien for Tiden har en ganske anden Karakter i
Tydskland end f. Ex. i England.

Dersom det da er sandt, hvad jo ofte gjentages,
at det er den klare sunde Forstand, forenet med en vis
Dybsindighed, men tillige med afgjort Uvillie mod mystiske
Spekulationer og ørkesløse Grublerier, der skal
afgive Maalestok for det danske Folks intellektuelle Begavelse
--- en Synsmaade, der finder Stadfæstelse i den
Kjendsgjerning, at Holberg, »Stifteren af vor Literatur,
han, hvis uberegnelige Indflydelse paa vor Smag og vort
naturlige Omdømme har gjort sig gjældende i næsten
alle Retninger«, netop er en af den sunde Menneskeforstands
genialeste Repræsentanter --- dersom dette, sige
vi, forholder sig saaledes: da maa der ogsaa kunne paavises
en til denne Aandsretning særlig svarende filosofisk
Opgave.

Opgaven er: \emph{at bestemme den menneskelige}
Videns Grændse.

Den første afgjørende Betingelse for at opnaae videnskabelig
Indsigt er, at jeg ved hvert Skridt nøiagtig
maa skjelne mellem det, jeg veed, og det, jeg ikke
veed. Mine Fremskridt i Viden beroe nu derpaa, at der
er Meget, som jeg endnu ikke veed, men som jeg dog
maaske med Tiden kan faae at vide. Der er i den
menneskelige Viden til alle Sider utallige Grændser, som
kunne overskrides. Ved nye Opdagelser, nye Synspunkter,
nye Methoder gjør Videnskaben uophørlig Fremskridt,
overskrider de gamle Grændser i alle Retninger,
og det vil i denne Forstand være umuligt at forudsige,
hvorvidt den i Tidernes Løb kan naae. Men Spørgsmaalet,
det for Liv og Videnskab afgjørende Grændse%
% --- p. 389 ---
\setcounter{footnote}{0}%
\opage{389}spørgsmaal, hvorpaa vi her i den sunde Menneskeforstands
Navn maae lægge Vægt, er et ganske andet.
Grundspørgsmaalet er, om Filosofien i Videns Form formaaer
at løse Tilværelsens Gaader.

Blandt Livggaaderne er der fornemmelig to, paa   % PRINTED AS IS (Bind 2, p. 389): „Livggaaderne“ for Livsgaaderne (zoom z4)
hvis Betydning den sunde Menneskeforstand altid har
lagt en særdeles Vægt, og for hvis Løsning den vistnok
vilde være Filosofien som Videnskab uendelig forbunden;
det er, som bekjendt, Realiteten af Guds Tilværelse
og Sjælens Udødelighed, hvorpaa den søger Bekræftelse.
Har Filosofien været istand til at give en
saadan Bekræftelse? Nei! Hvorfor ikke? Fordi den
sidste, afgjørende Besvarelse af de herhen hørende Spørgsmaal
ligge udenfor den menneskelige Videns Grændse.
At de mange hidtil opstillede Beviser, ved hvilke man
har søgt at grundfæste hine for Menneskeheden trøstelige
Sandheder, lade sig opløse i lige saa mange Tvivl,
har den negative Kritik tilstrækkelig godtgjort. Ved
at angribe Beviserne for de formentlige Sandheder er
man da gaaet til den modsatte Yderlighed; man har
saaledes ud fra den moderne Bevidstheds Standpunkt
opstillet Beviser, for hvis »ubønhørlige Logik« baade
Gud og Udødelighed maa forsvinde. Gud, det almægtig
villende Selv, er ikke til; et almægtigt, villende Væsen
er et begrebløst Begreb, en Troens Indbildning, en
Sæbeboble, der brister, naar den berøres af Kritikens
Naalespids: Guds Tilværelse er, som Tilværelsen af den
runde Firkant, umulig. Paa lignende Maade beviser
man, at Sjælens Udødelighed er umulig. Sjælens Liv
er oprindelig betinget af legemlige Organer; at tænke
uden Hjerne er ligesaa umuligt som at føle uden Nerver.
% --- p. 390 ---
\setcounter{footnote}{0}%
\opage{390}Altsaa: naar Nerverne døe, døer ogsaa Følelsen, naar
Legemet døer, døer ogsaa Sjælen.

Men hvad siger den sunde Menneskeforstand? Den
siger, at Filosofien er anmassende, naar den giver sig
Udseende af at vide det, som intet Menneske kan vide.
Den sunde Forstand tager ligesaa lidt Parti med den
Viden, der afsætter, som med den, der indsætter Gud;
ligesaa lidt med den, der modbeviser, som med den, der
beviser Sjælens Udødelighed. Den sunde Forstand er,
hvad de høieste Sandheder angaaer, en stor Skeptiker;
den kommer med sine Grunde for og imod aldrig ud over
Tvivlen. Men denne Forstandens uovervindelige Tvivlen
indeholder jo netop en Forudsætning, som Forstanden selv
maa ansee for hævet over enhver Tvivl, den Forudsætning,
at der i Virkeligheden er Noget, som vi i dette nærværende
Liv aldrig faae at vide, at den menneskelige
Viden, med andre Ord, har en uoverskridelig Grændse.

Der er mange Sætninger, der tage sig fornuftig ud,
saalænge de holdes i al Almindelighed; og dog, naar
der tales om noget Bestemt, og det gjælder om en virkelig
Anvendelse, kan det pludselig vise sig, at de ved
første Øiekast antagelige Sætninger indeholde en skjult
Tvetydighed og give en falsk Mening. Herhen høre den
abstrakte Filosofis almene Sætninger om den menneskelige
Videns Grændser. Skulde der, vil man sige, være
Noget for Viden aldeles Utilgængeligt, da maatte der
jo saavel i vort eget Indre som i Tingenes Verden være
Noget, der var aldeles utilgængeligt for Tanken, Noget
saa gaadefuldt, at det unddrog sig al vor Begriben, fordi
det unddrog sig alle saavel Naturens som Tankens Love.
Men et saadant hemmelighedsfuldt, i sig utænkeligt
% --- p. 391 ---
\setcounter{footnote}{0}%
\opage{391}Væsen er et sandt Uvæsen. Hvor farligt det er at give
sig hen til uklare Forestillinger om et mystisk Væsen,
hvis uforklarlige Virksomhed ikke er bunden til nogen
tænkelig Lov, hvis Aabenbarelser spotte alle Begreber
om Tingenes Sammenhæng ved betingede Aarsager og
tilsvarende Virkninger, kan sees af Fornuftens uafladelige
Kamp med Overtroen\footnote[1]{Overtroens Væsen er at antage Aarsager uden Love.}. I Kina f. Ex., hvor en forstandig
Opfattelse af de synlige Ting og deres Forbindelse
dog er naaet saavidt, at man til Hverdagsbehov er nogenlunde
hjemme i Astronomien, viser Almanaken udtrykkelig,
hvad en Kineser hver Dag Aaret rundt maa lide
af hemmelig Angst for Tilfældenes gaadefulde Spil; kloge
Forskrifter ere nødvendige. En kinesisk Kalender, begyndende
med den 8de Februar, Kinesernes Nyaarsdag, foreskriver:
»Den Første kan du gjøre dine religiøse Offringer
(saavel af Høns som Fisk); du kan henvende
dig til Himlen (i Taksigelser, Bønner, Løfter --- ved at
afbrænde Guldpapir, Straafigurer eller Fyrværkeri i uendelig
Mangfoldighed); du kan klæde dig i Galla, med
Skindhue eller elegant Livbaand\dots{}. Den Femte maa
du ikke begynde en Reise\dots{}. Den Ottende er tvivlsom\dots{}.
Den Tiende kan du indgaae Ægteskab\dots{}.
Den Trettende maa du Klokken fem om Morgenen sidde
med Ansigtet vendt mod Sydost«. Uden at miskjende den
europæiske Kulturs store Fortrin for den kinesiske, kan
man dog, naar man ret betænker, hvor farlig det endnu
i vor oplyste Tid er at være Tretten tilbords, at lægge
to Knive overkors, at tage en Knappenaal op om Mandagen,
begynde paa Noget om Fredagen o. s. v., med
% --- p. 392 ---
\setcounter{footnote}{0}%
\opage{392}Grund paastaae, at der midt i den civiliserede Verden
endda ikke er saa ganske Faa, som med sky Blik til
Tilfældets Dør, der altid staaer paa Klem, bære paa
en kinesisk Angst for hemmelige Magters lunefulde Indgriben
i Tingenes Sammenhæng.

Skal Slægten nogensinde naae saa vidt, at den omsider
faaer Bugt med al denne overtroiske Frygt og
bliver alle disse Ængstelser, disse uhyggelige Anelser
kvit: da maa den efterhaanden lære at gjøre sig fri
for mystiske Forestillinger om uopløselige Gaader. Men
hvad vil dette nu sige Andet, end at Tankens Lys i
Bevidstheden tilsidst maa blive saa klart, at alle Ting
i Verden vise sig som nødvendige Aabenbarelser af den
evige Fornuft. Fornuften i Verden og Fornuften i Mennesket
er eet og samme Væsen; Menneskets Viden om
Tingene svarer nøiagtig til dets Viden om sig selv. Er
det utænkeligt, at Noget kan være til, som ifølge sin
Natur falder udenfor Fornuftens Grændse, da er jo hermed
givet, at der umulig enten paa Jorden eller i Himlen
eller i Afgrundene kan være Noget til, som ifølge
sin Natur falder udenfor den menneskelige Videns Grændse.

Med den her opstillede Synsmaade gaaer det fortræffeligt,
saalænge man, som sagt, bevæger sig indenfor
de almindelige Betragtningers Omkreds. Ved idelig
at beraabe sig paa Fornuften og Lovene, Tænkningens
og Tingenes uforanderlige Love, kan man endog give
Sagen et Udseende, som om de strenge, forskende Videnskaber,
Fysik, Kemi, Fysiologi, Naturhistorie o, s. v.   % PRINTED AS IS (Bind 2, p. 392): „Naturhistorie o, s. v.“ comma for the stop of o. (b1-19-04)
alle forenede sig med Filosofien om at stadfæste den
Grundsætning, at den menneskelige Viden ikke har nogen
sidste endelig Grændse. Tage vi derimod Sagen som
% --- p. 393 ---
\setcounter{footnote}{0}%
\opage{393}den i Virkeligheden er, saa viser der sig noget Andet.
Thi saa viser det sig, at Grundsætningen om den menneskelige
Videns absolute Grændse ingenlunde staaer i
Strid med Grundsætningen om Lovenes Uforanderlighed;
saa viser det sig, at Naturvidenskaberne, langt fra at
nægte den skarpt betegnede Vidensgrændse, tvertimod,
alle for en og en for alle, forudsætte den og holde den
i Hævd.

Besinde vi os paa, hvorledes det egenlig er gaaet
til, at Naturvidenskaberne i den nyere Tid have kunnet
tage et saa mægtigt Opsving og opnaae saa store mærkværdige
Resultater, da er Grunden hertil bekjendt nok,
thi den ligger i Methoden. Methoden, den rette Fremgangsmaade,
bestemmes netop derved, at man strax fra
Begyndelsen af gjør sig Rede for, hvad der under givne
Betingelser kan undersøges og hvad der ligger udenfor
Undersøgelsen, hvad der kan vides og hvad der ikke
kan vides. Saalænge Grændsen mellem det, der kan
vides, og det, der ikke kan vides, endnu var svævende
og ubestemt, saaledes som i Oldtiden og gjennem hele
Middelalderen, saalænge var ogsaa Undersøgelsernes
Gang vaklende, og Udbyttet ubestemt; Stjernetyderi,
Guldmageri, Grublerier over de Vises Sten, slige mystiske
og magiske Bestræbelser vise tydelig nok, hvor
uklar Bevidstheden var over Videnskabens Grændser og
Opgaver.

Det er virkelig mærkværdigt at see den Skarphed,
hvormed de strenge Videnskaber udskille det, der
ligger udenfor Viden, fra det, som kan og maa blive
Gjenstand for videnskabelig Undersøgelse. »Formaalet
for Mekaniken«, hedder det, »er at udgrunde de almin%
% --- p. 394 ---
\setcounter{footnote}{0}%
\opage{394}delige Love for de mekaniske Naturkræfters Ligevægt
og for Massernes Bevægelser, hvorved en fuldkommen
Orden og Harmoni vedligeholdes i Legemverdenens Fænomener«.
Man skulde nu mene, at de mekaniske Kræfters
skjulte Væsen maatte fra Grunden af blive opklaret
i den analytiske Mekanik; men her gives udtrykkelig
den Erklæring, \emph{at vi ikke vide, hvad Kraft er.}
»Naar et Legeme« --- med disse Ord begynder Laplace
sin »Himlens Mekanik« --- »forandrer sin Stilling imod
et System af Legemer, hvilket vi ansee for hvilende,
saa synes det os at bevæge sig«; og han tilføier: »den
besynderlige Egenskabs Natur, som gjør, at et Legeme
forflyttes fra et Sted til et andet, er og vil bestandig
forblive ubekjendt; man har betegnet den med Navnet
Kraft, men man kan kun bestemme dens Virkninger og
de Love, hvorefter den virker«. En lignende Angivelse
af Videns Grændse møde vi i den kemiske Fysik, f. Ex.
i Læren om Varmen. »Følelsen er i Stand til ved Legemerne
at skjelne mellem forskjellige Tilstande, som
man udtrykker ved Ordene: \emph{hed}, \emph{varm}, \emph{kold} o. s. v.
Naar et Legeme, som man kalder koldt, bliver varmt
eller hedt, voxer dets Rumfang, det udvider sig. Den
\emph{ubekjendte Aarsag,} som bevirker Legemernes Udvidelse,
og som tillige frembringer de forskjellige, anførte
Indtryk paa vor Følelse, kalder man Varme«.
Altsaa: vi vide ikke, hvad Varme er; vi vide da heller
ikke, hvad Magnetisme er, vi vide ikke, hvad Elektricitet
er; vi vide ikke, hvad Aarsagen til en kemisk
Tiltrækning er; kort sagt: der er i Grundlaget for alle
naturvidenskabelige Lærebygninger nedlagt en Tilstaaelse
af, at den menneskelige Viden har en uoverskridelig
% --- p. 395 ---
\setcounter{footnote}{0}%
\opage{395}Grændse. Hermed have de exacte Videnskaber uden Omsvøb
indrømmet, at Tilværelsen trods Tingenes naturlige
Sammenhæng har sine Gaader; og dog er der vel
Ingen, der vil paastaae, at den forudsatte Vidensgrændse
vil indgyde Menneskene en hemmelig Angst for Naturens
dunkle Magter, eller bringe nogen Moders Sjæl til
at berede sig paa ubekjendte Muligheder ved at sidde
»den Trettende Klokken fem om Morgenen med Ansigtet
vendt mod Sydost«.

Betragte vi Forholdet mellem Sjæl og Legeme, da
kan Ingen tvivle om, at det i alle Henseender er bestemt
ved Love; ikke desto mindre maae Undersøgelserne
begynde fra to aldeles forskjellige Udgangspunkter.
Vælges Udgangspunktet i det Sjælelige, saa faae vi
Psykologien; vælges det i det Legemlige, faae vi Fysiologien.
Psykologien holder sig til de indre Erfaringer,
til vor umiddelbare Bevidsthed om Følelse, Fantasi,
Hukommelse, Forstand, Drift, Villie o. s. v. Fysiologien
holder sig til de ydre Erfaringer, den betragter
det menneskelige Legeme, som ethvert andet Legeme,
undersøger dets Organer og Væv, dets indre
Bygning og Sammensætning dels fra et mekanisk, dels
fra et kemisk, dels fra et morfologisk Standpunkt.

Nu er det vel saa, at Undersøgelsen af det indvortes
Sjælelige bestandig forudsætter Henblik til det
Legemlige, og at omvendt Udtydningen af de legemlige
Dannelser betinges af Henblik til de sjælelige Funktioner;
hvorledes skulde det vel være muligt for Psykologen
at gjøre sig noget Begreb om Synsvirksomheden
uden Kjendskab til Øiet, eller for Fysiologen at
tyde Øiets Bygning uden Hensyn til Synsvirksomheden?
% --- p. 396 ---
\setcounter{footnote}{0}%
\opage{396}Psykologien og Fysiologien maae komme hinanden imøde;
de maae flittig konferere og korrespondere med hinanden.

Desuagtet ere de to Videnskaber af en saa uligeartet
Karakter, at den ene udgjør en Del af Filosofien,
den anden af Fysiken. Der gives ingen ligefrem
Overgang fra den ene til den anden: Overgangen sker
ved et Spring. Følelse og Følenerve ere paa det Inderligste
forenede; men fra den fysiologiske Analyse af
en Følenerve til den psykologiske Erkjendelse af en Følelse,
hvilken uhyre Afstand! Gjenstanden for den udvortes
Erfaring er absolut uensartet med Gjenstanden
for den indvortes. Hvorlænge end Fysiologi og Psykologi
fortsætte deres Undersøgelser, og hvormeget de
end i Tidernes Løb kunne nærme sig til hinanden, der
kommer dog aldrig den Tid, da de tilfulde forstaae hinanden.
De tale hver sit Sprog, og de to Sprog ere
saa himmelvidt forskjellige, at hvad der siges i det ene,
ikke ved nogensomhelst Oversættelse lader sig gjengive i
det andet. Det er, som om Loven og Gaaden her havde
sat hinanden Stævne; Loven kan erkjendes, men Gaaden
er uopløselig. Naar en Følenerve overskæres, saa taber
den paagjældende Del af Legemet al Følelse; det er en
Lov, som Fysiologien har opdaget, og hvoraf den kan
gjøre en rig Anvendelse; men derved er den ikke kommen
Besvarelsen af det Spørgsmaal, hvorledes Nerven
forplanter Følelsen, et Haarsbred nærmere. Og hvad
Lighed er der saa mellem Hjernevindinger, Hjerneorganer,
den hvide og den graa Marv, og det Bevidsthedsliv
af Tanker og Forestillinger, der under Tilsyn af et styrende
Selv arbeider i Hjernen?

% --- p. 397 ---
\setcounter{footnote}{0}%
\opage{397}Det er aldeles umuligt, at enten en Hjernevirksomhed
skulde kunne frembringe en Tanke, eller en Tankevirksomhed
foraarsage en Bevægelse i Hjernen, dersom
der ikke er en indre fælles Natur, hvori de begge forenes.
En saadan Natur maa være legemlig som Hjernen
og sjælelig som Tanken; men en saadan Natur er
hverken Gjenstand for Tænkning eller for Sandsning;
den unddrager sig enhver mulig Undersøgelse: den ligger udenfor den menneskelige Videns Grændse.

Og dog maa der være en saadan Natur; vi vide
det med Vished: vi kunne bevise det psykologisk. Hvis
der i Hjernevirksomheden ikke var Andet end det, vi
ved Hjælp af fysiologisk Iagttagelse kunne forestille os:
da vilde en Hjernevirksomhed umulig kunne frembringe
en Tanke. Den fysiologiske Virksomhed er, hvad enten
den opfattes fra den mekaniske eller kemiske eller
morfologiske Side, legemlig, ene og alene legemlig. At
nu en Virksomhed i en legemlig Del kan frembringe
Forandringer i en anden legemlig Del, kunne vi vel forstaae,
thi her er Aarsagen ensartet med Virkningen;
men at en Virksomhed i det blot Legemlige kan frembringe
en Forandring i det aldeles Ulegemlige, er en
Modsigelse, thi her er Aarsagen absolut uensartet med
Virkningen. Hvorledes en Tanke kan frembringe en
anden Tanke, Tanken om Grund f. Ex. Tanken om
Følge, Tanken om Øiemed Tanken om Midler, er indlysende,
thi her er det Frembringende ensartet med det,
som frembringes; men at Tanken, en ulegemlig Aarsag,
skulde i Hjernen kunne frembringe en legemlig Virkning,
er utænkeligt, thi Aarsagen er igjen uensartet,
absolut uensartet med Virkningen.

% --- p. 398 ---
\setcounter{footnote}{0}%
\opage{398}Men, indvender man, naar vi nu vide, at Forbindelsen
mellem Sjæl og Legeme absolut maa beroe paa
en for begge fælles Natur, en Grundnatur, hvori det
Uligeartede er ligeartet, saa vide vi jo, at en saadan
Natur er til, og kunne altsaa ikke uden Modsigelse paastaae,
at dens Væren falder udenfor vor Viden. Det
Sofistiske i en saadan Indvending er iøinefaldende. Naar
vi sige, at der her er Noget, som ligger udenfor vor
Viden, saa sige vi ganske rigtig, at vi have en Viden
om, at et saadant Noget er, men vi sige jo ogsaa, at
vi ikke kunne komme ud over denne almene Viden og
bestemme det Ubekjendte nærmere; hermed er Grændsen
fastsat.

Sjælens lovbestemte Forbindelse med Legemet er
altsaa en Livsgaade. Saa enfoldig som denne Sætning
er, saa har den dog, hvad den sunde Menneskeforstand
let vil indsee, overordenlig vigtige Følger. Følgerne vise
sig strax, naar Talen er om Skilsmissen mellem Sjæl
og Legeme, om Døden; thi det forstaaer sig af sig selv:
som Forbindelsen er, maa ogsaa Skilsmissen være. Forholdet
mellem Loven og Gaaden bliver atter synligt.

Det er en Lov, ligesaa afgjørende for Psykologien
som for Fysiologien, at naar Tænkningsorganerne forstyrres,
saa forstyrres ogsaa Tænkningen, og at Bevidsthedens
Lys udslukkes, naar Legemet døer. Beraabende
sig paa denne Lov have Naturalisterne anset
det for afgjort, at Sjælen nødvendig maatte gaae tilgrunde
med Legemet. Det er en meget overilet Paastand.
Hverdagserfaringen lærer, at Bevidstheden kan
fordunkles, som f. Ex. i Søvne, ja endog forsvinde, som
% --- p. 399 ---
\setcounter{footnote}{0}%
\opage{399}i en dyb Besvimelse, og dog vende tilbage, fordi det
sjælelige Liv ikke ophører. Det Spørgsmaal: om Sjælen
er i Stand til at overleve Legemet, hænger paa det
Nøieste sammen med det Spørgsmaal: om det sjælelige
Liv trods sin Afhængighed af de legemlige Organer dog
ikke allerede i denne Tilværelse er hævet over Forkrænkeligheden
derved, at det i sit dybere Grundlag, i
den forborgne Natur, hvorpaa Foreningen af Sjæl og
Legeme oprindelig beroer, har Kræfter og Betingelser
for at danne sig et nyt Legeme, et Legeme af høiere
Orden, naar det jordiske er opløst.

»Mennesket døer som en Hund», siger en aandrig
Fysiolog; men, spørge vi, hvorledes døer saa Hunden?
Som et Menneske! Det er en skjøn Forklaring. Fysiologerne
selv tilstaae, at de ikke vide, hvad Livet er;
vi tage dem paa Ordet, og minde dem om, at de heller
ikke vide, hvad Døden er. Det er altsaa, thi her er
vor Videns Grændse, muligt, at Sjælen døer med Legemet,
men det er ogsaa --- thi Videns Grændse kan ikke
overskrides --- muligt, at Sjælen overlever Legemet.
Den Indvending, at et tilkommende Liv ikke lader sig
tænke, da ethvert Forsøg paa at danne sig en tydelig
Forestilling om en høiere Verdensorden opløser sig i
lutter Modsigelser, har ingen Gyldighed i den strenge
Videnskab. Fysiken, der holder sig strengt til de naturlige
Kjendsgjerninger, støder idelig paa Fænomener,
hvis indre Sammenhæng den ikke kan begribe. Den
kan, efter at have udfundet Lovene, lægge de anerkjendte
Love til Grund for en tilsvarende Forklaring,
og endda kan Forklaringen opløse sig i Modsi%
% --- p. 400 ---
\setcounter{footnote}{0}%
\opage{400}gelse\footnote[1]{Et Exempel. Ampère forklarer Magneten som bestaaende af mangfoldige molekulære, parallele og ensstillede elektriske Kredsløb. Af denne Hypothese kunne alle magnetiske Fænomener udtydes paa en til de mathematiske Formler nøiagtig svarende Maade. Men --- disse Kredsløb skulle nu foregaae midt i en Leder (Staalet f. Ex.), og maae da, ifølge Ampères egen Lov, forlade Molekület og udbrede sig over det hele Legeme. Tilstanden kan altsaa ikke vedblive, medmindre Molekulet bestandig frembringer et nyt Kredsløb. Men --- et elektrisk Kredsløb forudsætter et tilsvarende Forbrug og en tilsvarende Varmefrembringelse; men ved Magneten er der ingen Varmefrembringelse og intet Forbrug, altsaa heller intet elektrisk Kredsløb.}. Skulde Naturlæren fornægte de Fænomener, hvis
Væsen den ikke kan begribe, da maatte den ende med
at fornægte den hele Natur.

Paa den sikkre Fastsættelse af Videns Grændse
beroer Forholdet mellem Tro og Viden. Naar den yderste
Grændse for menneskelig Viden er naaet, vil Løsningen
af Livets høieste Opgaver altid stille sig saaledes,
at der fremkommer to lige gjældende og dog hinanden
ophævende Muligheder: det er muligt, at Sjælen
er dødelig, og det er muligt, at den er udødelig, det er
muligt, at der er en levende, personlig, villende Gud
til, men det Modsatte er ogsaa muligt. Først naar Videnskaben
efter alsidig Prøvelse har givet begrundet
Vished for, at den ene af de to Muligheder har samme
videnskabelige Gyldighed som den anden, har Viden sagt
sit sidste Ord og maa overlade Afgjørelsen til Troen.
Troen vælger nu den ene af de to Muligheder, Vantroen
den anden. Skilsmissen mellem Tro og Vantro
skeer ved Valget; Valget er ikke videnskabeligt, men
absolut personligt, her vælges i Kraft af personlige
Grunde.

% --- p. 401 ---
\setcounter{footnote}{0}%
\opage{401}Paa det ved Valget bestemte Grundlag kunne filosofiske
Lærebygninger atter opføres; men disse Lærebygninger
adskille sig nu fra de direkt videnskabelige
derved, at Grundene ere Personlighedsgrunde, at det
ikke er Nødvendigheden, men Friheden, hvorpaa Begrundelsen
hviler. Gaaden og Loven træde herved ind
i en ny Fase.

Det er en ligesaa bekjendt som forslidt Indvending,
at Personlighedslæren ved i Frihedens Navn at gjøre
Troen til selvstændigt Erkjendelsesprincip indfører Vilkaarligheden,
gjør Brud paa Lovene og kommer i Strid
med Tingenes naturnødvendige Sammenhæng. Navnlig
er det Læren om Underet, der vækker Forargelse. Det
kan, siger man, fra Videnskabens Standpunkt vel indrømmes,
at der er meget Gaadefuldt i Naturen. Naturlæren
opdager jo bestandig nye Fænomener, nye
Kræfter og nye Muligheder; men disse Opdagelser kunne
ingenlunde svække vor Overbevisning om Tingenes nødvendige
Sammenhæng, da enhver ny Opdagelse er, om
ikke altid en Opdagelse af nye Love, saa dog en Udvidelse
og Stadfæstelse af de allerede bekjendte. Med
Villiesprincipet indføres derimod Vilkaarligheden; men
Vilkaarligheden, de vilde Muligheders Løsen, er uforenelig
med Lovene.

Hvorledes det forholder sig med denne Sag, kan
oplyses saavel fra en menneskelig som fra en guddommelig
Side. Vi tillægge jo ikke blot Menneskene,
men ogsaa Dyrene, vilkaarlige Bevægelser, Bevægelser
der have deres Aarsag i umiddelbar Selvbestemmelse.
Dersom det nu virkelig var sandt, at enhver Yttring af
vilkaarlig Selvbestemmelse stred mod Naturlovene, da
% --- p. 402 ---
\setcounter{footnote}{0}%
\opage{402}maatte man nødvendig antage Eet af To: enten er Vilkaarligheden
et Skin, og der er slet intet Vilkaarligt,
eller, hvis der gives vilkaarlige Livsyttringer, saa maa
Vilkaarligheden hvert Øieblik gjøre Brud paa Lovene.
Maae nu begge Antagelser erklæres for falske, saa er
hermed tillige givet, at selvbestemmende Væseners Vilkaarlighed
meget vel kan forenes med uforanderlige
Love.

Der er i Tilværelsen en gjennemgaaende Modsætning
mellem Naturens og Aandens, Nødvendighedens og
Frihedens Væsen, men der er i Principet selv ligesaa lidt
nogen Strid imellem Lovene for de to Regioner, som
der indenfor Naturrigets Omkreds er nogen Strid mellem
de mekaniske og kemiske Love.

Ja, vedbliver man fra Kritikens Side, Sligt lader sig
nok saadan høre i Almindelighed, og saalænge det kun høres
i al Almindelighed, klinger det ret godt; men drage vi
Tilfældene frem, saa viser sig strax det Urimelige. At
den menneskelige Villies Vilkaarlighed ikke kan bryde
Naturens Sammenhæng, følger ligefrem af, at Menneskets
Villie er afmægtig; men ifølge Personlighedslærens
øverste Grundtanke, skal Guds almægtige Villie
jo kunne gribe vilkaarlig ind i Naturens Orden, gjøre
Mirakler og altsaa bryde Lovene. --- Her er nu atter
en Blanding af Sandt og Falsk. Det er ganske vist et
Særmærke for Underet, at det har den almægtige Villie
til frit bestemmende Aarsag; men det er fra Troens og
Personlighedens Synspunkt ingenlunde Underets Kjende,
at det skulde medføre Brud paa de guddommelige
Love. At gjøre det Lovstridige til Underets Kjendemærke
er et Paafund af den nyere Kritik. Grunden til,
% --- p. 403 ---
\setcounter{footnote}{0}%
\opage{403}at dette Paafund har vundet saa stærk Indgang, at Sætningen
nu er nær ved at blive ophøiet til en af Videnskabens
store Grundsætninger, ligger klart for Dagen.
For at tage sig ud som Videnskab og fuske i at forklare
eller dog forsvare Miraklet, har Theologien fra
Arilds Tid af maattet overskride Videns Grændse; for
at gjendrive den fuskeragtige Forklaring og blive Miraklet
kvit, er den negative Kritik fra sin Side ligeledes
bleven fristet til at overskride Videns Grændse. Theologien
slutter: naar Gud er almægtig, saa maa han
efter Behag kunne tilintetgjøre de naturlige Aarsager
og suspendere Lovene; Kritiken slutter: naar Gud gjør
en Gjerning, der ikke lader sig forklare efter de os bekjendte
Love, saa gjør han en Ugjerning og kommer i
Modsigelse med sig selv. Det er falsk Logik\footnote[1]{Ved at forandre Udtrykket og sige: en Begivenhed, der strider mod en eneste af de os bekjendte Love, strider imod alle Love, vinder man kun en tvetydig Fordel. Udtrykket: at stride mod en Lov, kan nemlig tages i to Betydninger. Det kan betyde: at staae i Modsigelse med Lovens Begreb som Lov (at tegne en rund Firkant, fremvise en Virkning uden Aarsag, sætte 2 + 2 = 5 o. desl.). I denne Mening finder Udtrykket ingen Anvendelse paa Miraklet. Det kan dernæst betyde: at være i tilsyneladende Strid med en virkelig Lov, saa at Opfyldelsen af den ubekjendte Lov synes at ophæve den bekjendte. I denne Mening finder Udtrykket Anvendelse ikke alene paa Miraklet, men ogsaa paa Naturbegivenhederne indbyrdes. Forestille vi os, at En, der hidtil kun havde lært at forklare Fænomener efter mekaniske Love, pludselig blev Vidne til overraskende Fænomener af magnetiske, elektriske, elektromagnetiske Love: mon han da ikke ville mene at see lutter Vidundere?}. Den
tomme Ræsonneren over, hvad Gud kan gjøre, og hvad
Gud ikke kan gjøre, forvirrer Hjernerne.

Ifølge det her Anførte vil nu Enhver med sund
Forstand og almindelig Dannelse kunne sige sig selv,
% --- p. 404 ---
\setcounter{footnote}{0}%
\opage{404}hvor vigtigt det er, at der een Gang for alle sættes en
Grændse mellem det, som kan vides, og det, som ikke
kan vides; først naar Grændsen er sat, og dens Gyldighed
anerkjendt, ville Liv og Videnskab kunne mødes
i ærlig Forstaaelse og see hinanden i Øiet. En Opgave
af saa gjennemgribende Betydning vil naturligvis
efterhaanden trænge sig frem og blive behandlet i alle
Lande af alle Tænkere, saa langt som Videnskaben
naaer; men det udelukker ikke, at dens Løsning kan
ligge den ene Folkeindividualitet nærmere end den anden.
Den tydske Spekulation er gaaet ud over Grændsen;
den engelske Erfaringsfilosofi er ikke naaet til Grændsen;
den franske Forstandslogik med de exakte og dog
bevægelige Former gjør Grændsen flydende abstrakt ved
at gjøre den mathematisk; det er da ikke urimeligt, at
et Folk, hvis eiendommelige Literatur har en Holberg
til Fader og en dybsindig Mythologi til Moder, et Folk,
med hvis jævne sunde Forstand det saa naturlig stemmer,
at det ikke har forhastet sig med at bygge himmelstormende
Systemer, men givet sig Tid og først begyndt
at filosofere paa egen Haand, efter at det i Naturlære,
Historie og Poesi har sikkret sig et forsvarligt
Grundlag for Aandslæren, at, sige vi, et saadant Folk
maaske kunde faae Øiet op for en af Nutidens væsenligste
Opgaver, at dets sildig fødte Filosofi kunde voxe
ved Opgaven og fremme dens Løsning ved at yde, om
end kun et Bidrag; et Bidrag er altid Noget.
\signature{R. \emph{Nielsen}.}

\end{document}
